% Editable publishing source. Rebuild without --refresh-tex to preserve edits.
\documentclass[10pt,a4paper]{article}
% Shared editable typography for all five document editions.
\usepackage{fontspec}
\setmainfont{DejaVu Serif}
\setsansfont{DejaVu Sans}
\setmonofont{DejaVu Sans Mono}[Scale=.86]
\usepackage{polyglossia}
\setdefaultlanguage{russian}
\setotherlanguage{english}
\usepackage[a4paper,margin=18mm,top=21mm,bottom=21mm,headheight=15pt]{geometry}
\usepackage{xcolor}
\definecolor{Accent}{HTML}{185A86}
\definecolor{Muted}{HTML}{526573}
\usepackage{graphicx}
\usepackage{longtable,booktabs,array,calc}
\usepackage{fvextra}
\usepackage{pdflscape}
\usepackage{titlesec}
\usepackage{fancyhdr}
\usepackage{enumitem}
\usepackage{etoolbox}
\usepackage{needspace}
\usepackage{tocloft}
\usepackage[unicode,colorlinks=true,linkcolor=Accent,urlcolor=Accent,bookmarksnumbered=false]{hyperref}
\usepackage{bookmark}
\usepackage{microtype}
\AtBeginDocument{\raggedright}
\setcounter{secnumdepth}{0}
\setcounter{tocdepth}{1}
\setlength{\parindent}{0pt}
\setlength{\parskip}{4pt plus 1pt}
\linespread{1.0}
\setlength{\emergencystretch}{3em}
\setlength{\tabcolsep}{4pt}
\renewcommand{\arraystretch}{1.16}
\AtBeginEnvironment{longtable}{\small}
% Measure a compact table or its header and first two data rows. The output
% remains a breakable longtable. Cap the reservation for unusually tall rows.
\newsavebox{\DocumentTableStart}
\newlength{\DocumentTableSpace}
\newcommand{\KeepTableStart}[2]{%
  \setbox\DocumentTableStart=\vbox{\hsize=\linewidth\small #2}%
  \setlength{\DocumentTableSpace}{\dimexpr\ht\DocumentTableStart+\dp\DocumentTableStart
    +\LTpre+\baselineskip+#1\baselineskip*4\relax}%
  \ifdim\DocumentTableSpace>.5\textheight
    \setlength{\DocumentTableSpace}{.5\textheight}%
  \fi
  \Needspace{\DocumentTableSpace}%
}
% Measure short lists once. Small blocks move together; long ones are unboxed
% into the normal page builder and remain breakable, never scaled or clipped.
\newsavebox{\DocumentShortBlock}
\newcommand{\KeepShortBlock}[1]{%
  % A heading's page-space check must not run inside the measuring box: it
  % would inspect the outer page and embed a spurious forced break in the box.
  \par\setbox\DocumentShortBlock=\vbox{\hsize=\linewidth
    \renewcommand{\Needspace}[1]{}#1\par}%
  \ifdim\dimexpr\ht\DocumentShortBlock+\dp\DocumentShortBlock\relax<.32\textheight
    \Needspace{\dimexpr\ht\DocumentShortBlock+\dp\DocumentShortBlock+\medskipamount\relax}%
  \else
    \Needspace{5\baselineskip}%
  \fi
  % Unboxing does not add normal inter-paragraph baseline glue at the leading
  % edge. Give the measured block standard vertical separation from prior text.
  \addvspace{\medskipamount}%
  \unvbox\DocumentShortBlock
}
\preto\section{\Needspace{5\baselineskip}}
\preto\subsection{\Needspace{5\baselineskip}}
\preto\subsubsection{\Needspace{4\baselineskip}}
\clubpenalty=10000
\widowpenalty=10000
\setlist{leftmargin=1.5em,itemsep=2pt,topsep=3pt}
\providecommand{\tightlist}{\setlength{\itemsep}{1pt}\setlength{\parskip}{0pt}}
\titleformat{\section}{\large\sffamily\bfseries\color{Accent}}{}{0pt}{}
\titleformat{\subsection}{\normalsize\sffamily\bfseries\color{Accent}}{}{0pt}{}
\titleformat{\subsubsection}{\normalsize\sffamily\bfseries\color{Accent}}{}{0pt}{}
\titlespacing*{\section}{0pt}{18pt}{8pt}
\titlespacing*{\subsection}{0pt}{13pt}{6pt}
\titlespacing*{\subsubsection}{0pt}{10pt}{5pt}
\pagestyle{fancy}
\fancyhf{}
\fancyfoot[C]{\small\color{Muted}\thepage}
\renewcommand{\headrulewidth}{.35pt}
\renewcommand{\footrulewidth}{0pt}
% A single compact opening page, with profile-controlled TOC depth. Reserve a
% wider page-number column for the multi-page BRD and example editions.
\setlength{\cftbeforesecskip}{5pt}
\setlength{\cftbeforesubsecskip}{2pt}
\renewcommand{\cftsecfont}{\small\sffamily\bfseries}
\renewcommand{\cftsubsecfont}{\small\sffamily}
\renewcommand{\cftsecpagefont}{\small}
\renewcommand{\cftsubsecpagefont}{\small}
\cftsetpnumwidth{2em}
\cftsetrmarg{3em}
\newcommand{\DocumentTitle}[3]{%
  \thispagestyle{plain}
  {\sffamily\bfseries\color{Accent}\fontsize{23}{29}\selectfont #1\par}
  \vspace{3mm}
  {\sffamily\small\color{Muted}#2\par}
  \ifstrempty{#3}{}{\vspace{4mm}#3\par}%
  \begingroup
  \renewcommand{\contentsname}{Содержание}%
  \tableofcontents
  \endgroup
  \vspace{5mm}
  {\color{Accent}\rule{\linewidth}{.8pt}\par}
  \clearpage
}
% Compact standalone profile: keep the title and body on the same opening page.
% It intentionally has no table of contents or forced title-page break.
\newcommand{\CompactDocumentTitle}[3]{%
  \thispagestyle{plain}
  {\sffamily\bfseries\color{Accent}\fontsize{23}{29}\selectfont #1\par}
  \vspace{3mm}
  {\sffamily\small\color{Muted}#2\par}
  \ifstrempty{#3}{}{\vspace{4mm}#3\par}%
  \vspace{5mm}
  {\color{Accent}\rule{\linewidth}{.8pt}\par}
  \vspace{4mm}
}

\hypersetup{pdftitle={SolutionArchitectAI},pdfauthor={},pdfsubject={Бизнес-требования · MVP}}
\fancyhead[L]{\small SolutionArchitectAI}
\fancyhead[R]{\small Бизнес-требования · MVP}
\begin{document}
\DocumentTitle{SolutionArchitectAI}{Бизнес-требования · MVP}{Назначение, этапы проектирования, результаты и критерии приёмки первой версии.}
\section{Цель и пользователи}\label{ux446ux435ux43bux44c-ux438-ux43fux43eux43bux44cux437ux43eux432ux430ux442ux435ux43bux438}

Помочь solution-архитектору преобразовать бизнес-требования в проверяемое описание потоков данных и межсистемных взаимодействий: уточнить пробелы, получить предложения ИИ, внести правки и утвердить результат.

Основной пользователь --- архитектор. Бизнес-аналитик помогает уточнять требования. Окончательные архитектурные решения принимает человек.

\section{Входные данные}\label{ux432ux445ux43eux434ux43dux44bux435-ux434ux430ux43dux43dux44bux435}

\textbf{Бизнес-требования, BRD:}

\begin{itemize}
\tightlist
\item
  назначение проектируемой системы;
\item
  её пользователи и другие участники бизнес-процессов;
\item
  функциональные требования;
\item
  нефункциональные требования;
\item
  используемые данные;
\item
  ограничения.
\end{itemize}

\textbf{Корпоративный контекст, если предоставлен:}

\begin{itemize}
\tightlist
\item
  существующие системы, интерфейсы и активы;
\item
  доступные технологии и архитектурные паттерны;
\item
  архитектурные стандарты и требования безопасности;
\item
  правила, условия их применения и допустимые исключения.
\end{itemize}

\textbf{Ответы и решения архитектора:}

\begin{itemize}
\tightlist
\item
  ответы на вопросы ИИ;
\item
  принятые допущения;
\item
  правки и согласования.
\end{itemize}

Требования и корпоративный контекст могут поступать из документов, структурированных данных и внешних источников. Их смысл не привязан к конкретному файловому формату или хранилищу. Перечень поддерживаемых способов получения данных для первой версии определяется при техническом проектировании.

Без корпоративного контекста работа допустима: неизвестные сведения остаются вопросами, а выбранные ИИ технологии и классификации --- предложениями. Недостаток информации не заполняется неподтверждёнными фактами.

\section{Этапы проектирования}\label{ux44dux442ux430ux43fux44b-ux43fux440ux43eux435ux43aux442ux438ux440ux43eux432ux430ux43dux438ux44f}

Основные этапы: \textbf{анализ требований → архитектура данных → межсистемные взаимодействия}. После каждого этапа архитектор проверяет и подтверждает результат. Из принятых результатов формируется итоговый комплект артефактов.

\subsection{FR-01. Анализ требований}\label{fr-01-ux430ux43dux430ux43bux438ux437-ux442ux440ux435ux431ux43eux432ux430ux43dux438ux439}

ИИ выделяет:

\begin{itemize}
\tightlist
\item
  цели и назначение проектируемой системы;
\item
  границы её ответственности: какие функции входят в систему, какие остаются внешним системам или участникам;
\item
  категории пользователей системы и их бизнес-роли;
\item
  функциональные и нефункциональные требования;
\item
  источники, состав и назначение данных;
\item
  периодичность и условия обмена данными;
\item
  архитектурно значимые ограничения.
\end{itemize}

Для извлечённого факта сохраняются подтверждающий фрагмент и его место во входном документе или другом источнике.

\subsection{FR-02. Архитектура данных}\label{fr-02-ux430ux440ux445ux438ux442ux435ux43aux442ux443ux440ux430-ux434ux430ux43dux43dux44bux445}

ИИ предлагает источники и наборы данных, их классификацию, хранение и потоки. Для потока указываются отправитель, получатель, состав данных, назначение и периодичность передачи.

\subsection{FR-03. Межсистемные взаимодействия}\label{fr-03-ux43cux435ux436ux441ux438ux441ux442ux435ux43cux43dux44bux435-ux432ux437ux430ux438ux43cux43eux434ux435ux439ux441ux442ux432ux438ux44f}

ИИ предлагает системы и логические сервисы, интерфейсы, инициаторов и получателей вызовов, протоколы, режимы обмена и свойства защиты. Технические взаимодействия связываются с потоками данных. Для передачи результата внешнему получателю указываются класс информации, разрешённый получатель и соответствующий правилам защищённый канал.

При выборе границ сервисов учитывается влияние загрузки данных на работу пользователей. Если данные поступают в пользовательское время, разделение загрузки и пользовательского обслуживания рассматривается с учётом требований и правил контекста.

\section{Сквозные возможности}\label{ux441ux43aux432ux43eux437ux43dux44bux435-ux432ux43eux437ux43cux43eux436ux43dux43eux441ux442ux438}

Уточнения, подбор и обоснование решений, правки человека и проверки выполняются на протяжении всего процесса. Номера функций служат для ссылок на требования и не задают порядок их выполнения.

\subsection{FR-04. Вопросы ИИ и ответы человека}\label{fr-04-ux432ux43eux43fux440ux43eux441ux44b-ux438ux438-ux438-ux43eux442ux432ux435ux442ux44b-ux447ux435ux43bux43eux432ux435ux43aux430}

При неполных или противоречивых вводных ИИ формулирует вопросы архитектору или бизнес-аналитику. Пользователь видит, что требуется уточнить, почему это важно и какое решение зависит от ответа.

Если вопрос блокирующий, система ожидает ответа либо явно принятого допущения и не продолжает зависимое проектирование. Ответ сохраняется и учитывается при уточнении требований и модели. Неблокирующие вопросы остаются видны в результатах. Такой цикл возможен на любом этапе.

\subsection{FR-05. Подбор и обоснование решений}\label{fr-05-ux43fux43eux434ux431ux43eux440-ux438-ux43eux431ux43eux441ux43dux43eux432ux430ux43dux438ux435-ux440ux435ux448ux435ux43dux438ux439}

При проектировании данных и взаимодействий ИИ сопоставляет требования с доступными активами, стандартами и подходящими архитектурными паттернами, рассматривает применимые варианты и предлагает решение \textbf{вместе с его обоснованием}.

Предложение содержит основания выбора, затронутые требования, применённые правила или паттерны, существенную альтернативу и статус решения. Если основания недостаточны, ИИ задаёт вопрос или явно обозначает допущение. Обоснование доступно до принятия решения человеком и обновляется при его изменении.

\subsection{FR-06. Участие архитектора на всём протяжении работы}\label{fr-06-ux443ux447ux430ux441ux442ux438ux435-ux430ux440ux445ux438ux442ux435ux43aux442ux43eux440ux430-ux43dux430-ux432ux441ux451ux43c-ux43fux440ux43eux442ux44fux436ux435ux43dux438ux438-ux440ux430ux431ux43eux442ux44b}

Архитектор может на любом этапе и после его завершения задавать уточнения, отвечать на вопросы ИИ, исправлять требования и модель, изменять, принимать или отклонять решения. Участие человека не ограничивается отдельной стадией согласования.

При добавлении объекта через ИИ служебный идентификатор назначает система. Человеку показываются понятные названия; вручную придумывать и запоминать коды объектов не требуется. После прямой правки YAML проверяется целостность ссылок.

Подтверждённые решения сохраняются при последующей обработке. Изменение вводных или принятого решения создаёт новую ревизию; согласования изменённых и зависимых результатов теряют актуальность. В MVP достаточно повторно обработать последующие результаты целиком.

\subsection{FR-07. Проверки и понятный отчёт}\label{fr-07-ux43fux440ux43eux432ux435ux440ux43aux438-ux438-ux43fux43eux43dux44fux442ux43dux44bux439-ux43eux442ux447ux451ux442}

Авторитетные источники, на основании которых проводится проверка:

\begin{itemize}
\tightlist
\item
  схема модели и её обязательные условия: уникальность идентификаторов, существование ссылок, заполнение обязательных сведений;
\item
  заданные архитектурные стандарты, требования безопасности и правила допустимых элементов и связей выбранной нотации;
\item
  требования проекта, ответы человека, принятые решения и предусмотренные правилами исключения.
\end{itemize}

Для применённого правила в отчёте сохраняется ссылка на источник и его редакцию.

Формализованные условия и бизнес-правила проверяются детерминированно. ИИ и эвристики могут выявлять возможные смысловые противоречия и пробелы; такие замечания показываются как предположения или вопросы, требующие проверки человеком. Для неформализованного требования нельзя заявлять успешную автоматическую проверку.

Пользователь получает отчёт: что проверено, по какому правилу или требованию, какой объект затронут, что нарушено и что требуется уточнить. Различимы результаты «пройдено», «нарушение», «не применимо» и «не проверено», а также формальные ошибки и предположительные замечания ИИ. Проверки повторяются после изменений; выбор алгоритмов и движка правил относится к техническому проектированию.

Блокирующие формальные нарушения требуют исправления. Согласованием нельзя скрыть ошибку; исключение допустимо, только если предусмотрено самим правилом.

\section{Правила работы}\label{ux43fux440ux430ux432ux438ux43bux430-ux440ux430ux431ux43eux442ux44b}

\begin{itemize}
\tightlist
\item
  Факт, ответ человека, допущение и предложение ИИ различимы. Совпадение цитаты с документом подтверждает происхождение текста, но не правильность его трактовки.
\item
  Направление движения данных и направление технического вызова различаются: сервис может вызвать источник и получить от него данные.
\item
  Неизвестные протоколы, порты, средства защиты и характеристики систем не выдаются за установленные. Если сведения нужны для согласования, они уточняются.
\item
  Общие стандарты и требования безопасности подключаются из авторитетных источников. Изменение проектного решения в диалоге не меняет корпоративное правило и не создаёт неразрешённое исключение.
\end{itemize}

\section{Выходные артефакты}\label{ux432ux44bux445ux43eux434ux43dux44bux435-ux430ux440ux442ux435ux444ux430ux43aux442ux44b}

\subsection{FR-08. Формирование результатов}\label{fr-08-ux444ux43eux440ux43cux438ux440ux43eux432ux430ux43dux438ux435-ux440ux435ux437ux443ux43bux44cux442ux430ux442ux43eux432}

Из одной принятой модели формируются согласованные таблицы, две диаграммы, архитектурный манифест и отчёт проверок. Манифест предоставляется в YAML для чтения и правок человеком.

\begin{longtable}[]{@{}
  >{\raggedright\arraybackslash}p{(\columnwidth - 4\tabcolsep) * \real{0.1400}}
  >{\raggedright\arraybackslash}p{(\columnwidth - 4\tabcolsep) * \real{0.2800}}
  >{\raggedright\arraybackslash}p{(\columnwidth - 4\tabcolsep) * \real{0.5800}}@{}}
\toprule\noalign{}
\begin{minipage}[b]{\linewidth}\raggedright
ID
\end{minipage} & \begin{minipage}[b]{\linewidth}\raggedright
Артефакт
\end{minipage} & \begin{minipage}[b]{\linewidth}\raggedright
Содержание
\end{minipage} \\
\midrule\noalign{}
\endhead
\bottomrule\noalign{}
\endlastfoot
A-01 & Реестр требований и уточнений & Требования с происхождением, вопросы, ответы, допущения и статусы \\
A-02 & Каталог данных и матрица потоков & Источники, наборы данных, классификация, направления и периодичность \\
A-03 & Диаграмма потоков данных & Перемещение данных между участниками и системами в согласованной нотации ArchiMate \\
A-04 & Каталог систем и матрица интеграций & Системы, сервисы, интерфейсы и технические взаимодействия, связанные с потоками \\
A-05 & Диаграмма межсистемных взаимодействий & Техническое представление интерфейсов, вызовов и протоколов \\
A-06 & Архитектурный манифест & Машиночитаемое описание контекста, систем, данных, потоков, интерфейсов и свойств защиты для дальнейшего проектирования. Реестр требований и согласования ведутся отдельно \\
A-07 & Реестр решений и отчёт проверок & Обоснования, сопровождающие предложения на каждом этапе, альтернативы, результаты проверок и оставшиеся ограничения \\
\end{longtable}

Объекты имеют устойчивые служебные идентификаторы. В материалах для человека используются понятные названия систем и данных; ссылки в YAML задаются осмысленными ключами. Таблицы и диаграммы представляют одну редакцию модели. Диаграммы доступны для просмотра и повторного формирования; манифест --- для редактирования и дальнейшей обработки.

\subsection{Когда создаются артефакты}\label{ux43aux43eux433ux434ux430-ux441ux43eux437ux434ux430ux44eux442ux441ux44f-ux430ux440ux442ux435ux444ux430ux43aux442ux44b}

\begin{longtable}[]{@{}
  >{\raggedright\arraybackslash}p{(\columnwidth - 4\tabcolsep) * \real{0.3100}}
  >{\raggedright\arraybackslash}p{(\columnwidth - 4\tabcolsep) * \real{0.3500}}
  >{\raggedright\arraybackslash}p{(\columnwidth - 4\tabcolsep) * \real{0.3400}}@{}}
\toprule\noalign{}
\begin{minipage}[b]{\linewidth}\raggedright
Этап
\end{minipage} & \begin{minipage}[b]{\linewidth}\raggedright
Создаются
\end{minipage} & \begin{minipage}[b]{\linewidth}\raggedright
Дополняются
\end{minipage} \\
\midrule\noalign{}
\endhead
\bottomrule\noalign{}
\endlastfoot
Анализ требований, FR-01 & A-01 --- реестр требований; начальные A-06 --- манифест и A-07 --- отчёт проверок & При уточнениях обновляются требования и их происхождение \\
Архитектура данных, FR-02 & A-02 --- каталог и матрица потоков; A-03 --- диаграмма потоков & A-01, A-06, A-07, включая обоснование предлагаемых решений \\
Межсистемные взаимодействия, FR-03 & A-04 --- каталог систем и матрица интеграций; A-05 --- диаграмма взаимодействий & A-01, A-06, A-07; при изменении потоков также A-02 и A-03 \\
\end{longtable}

Решения и проверки в A-07 накапливаются по мере работы. FR-08 формирует полный комплект из последних согласованных редакций; это не отдельный этап принятия архитектурных решений.

\section{Ограничения и тестирование}\label{ux43eux433ux440ux430ux43dux438ux447ux435ux43dux438ux44f-ux438-ux442ux435ux441ux442ux438ux440ux43eux432ux430ux43dux438ux435}

Разработка и приёмка выполняются \textbf{на синтетических бизнес-требованиях и синтетических корпоративных данных}: каталогах систем и активов, интерфейсах, стандартах, политиках и исключениях. Закрытая инфраструктура для проверки MVP не требуется.

Для проверки используется ручной эталон: входные требования и контекст, принятые допущения и полный ожидаемый комплект A-01--A-07. Проверочные сценарии включают неполные вводные, противоречащие правила, изменение периодичности обмена, класса данных и требований к доступности.

Контекст и правила имеют версии. Для принятой модели повторное формирование представлений без ИИ сохраняет смысловое содержимое. Повторная генерация ИИ может дать иной черновик, но не должна молча заменять подтверждённые решения.

Неудачная обработка явно отмечается; неполный результат не объявляется согласованным. Использование реальных документов требует определённых правил защиты и анонимизации; тестовый заменитель такого механизма не подтверждает безопасность реальных данных.

\section{Критерии приёмки}\label{ux43aux440ux438ux442ux435ux440ux438ux438-ux43fux440ux438ux451ux43cux43aux438}

\begin{longtable}[]{@{}
  >{\raggedright\arraybackslash}p{(\columnwidth - 4\tabcolsep) * \real{0.1400}}
  >{\raggedright\arraybackslash}p{(\columnwidth - 4\tabcolsep) * \real{0.2800}}
  >{\raggedright\arraybackslash}p{(\columnwidth - 4\tabcolsep) * \real{0.5800}}@{}}
\toprule\noalign{}
\begin{minipage}[b]{\linewidth}\raggedright
ID
\end{minipage} & \begin{minipage}[b]{\linewidth}\raggedright
Проверка
\end{minipage} & \begin{minipage}[b]{\linewidth}\raggedright
Признак приёмки
\end{minipage} \\
\midrule\noalign{}
\endhead
\bottomrule\noalign{}
\endlastfoot
AC-01 & Уточнение требований & ИИ задаёт вопросы о блокирующих пробелах согласованного набора примеров; до ответа или принятого допущения зависимое проектирование не продолжается; полученный ответ учитывается в результате \\
AC-02 & Происхождение и обоснование & Факты связаны с проверяемыми входными сведениями; добавленные решения имеют статус и основание \\
AC-03 & Сквозной комплект & Получены A-01--A-07, обе диаграммы читаемы, объекты и связи согласованы; формальных ошибок нет \\
AC-04 & Правки на разных этапах & Пользователь может изменить требования, поток или интерфейс в процессе проектирования и после согласования; правки сохраняются при повторной обработке \\
AC-05 & Новая ревизия & Изменение вводных снимает актуальность зависимых согласований \\
AC-06 & Нарушения & Выявлены повторный идентификатор, отсутствующая ссылка и нарушение заданного правила; невыполненная проверка видна явно \\
AC-07 & Корпоративный контекст & На синтетических примерах проверены работа с контекстом и без него, применение ограничений и предусмотренных правилом исключений \\
AC-08 & Обоснование до согласования & Предложение по данным или взаимодействиям сразу содержит основание выбора; оно доступно человеку до принятия решения и обновляется при изменении предложения \\
AC-09 & Прозрачность проверок & Для замечания видны объект и основание; формальное нарушение, предположение ИИ и невыполненная проверка различимы \\
\end{longtable}

Сравнивается смысл результата: обязательные объекты, связи, решения и соблюдение правил. Дословное совпадение текста ИИ и расположение блоков на диаграмме не требуются.

\section{За пределами MVP}\label{ux437ux430-ux43fux440ux435ux434ux435ux43bux430ux43cux438-mvp}

Физическая модель БД и ER; Sequence-диаграммы; полная спецификация API; сетевая топология и назначение IP; конфигурации развёртывания; генерация политик доступа; автоматическое сопровождение существующих архитектур; запись в корпоративные реестры; визуальный редактор и многопользовательская работа.

Требования доступности, восстановления и безопасности анализируются как ограничения, даже если расчёт их технической реализации отложен. Эксплуатационное подключение к корпоративным платформам и двусторонняя синхронизация моделей требуют отдельного расширения объёма.

\end{document}
